% Generated draft: requires mathematical and source review.
\documentclass[11pt]{article}
\usepackage[margin=25mm]{geometry}
\usepackage{fontspec}
\setmainfont{Latin Modern Roman}
\usepackage{amsmath,amssymb,mathtools,unicode-math}
\setmathfont{Latin Modern Math}
\usepackage{enumitem,needspace,xurl,hyperref}
\setlength{\parindent}{0pt}
\setlength{\parskip}{5pt}
\setlength{\emergencystretch}{3em}
\newcommand{\N}{\mathbb{N}}
\newcommand{\Z}{\mathbb{Z}}
\newcommand{\Q}{\mathbb{Q}}
\newcommand{\R}{\mathbb{R}}
\newcommand{\C}{\mathbb{C}}
\newcommand{\F}{\mathbb{F}}
\newcommand{\A}{\mathbb{A}}
\newcommand{\PP}{\mathbb{P}}
\newcommand{\E}{\mathbb{E}}
\newcommand{\eps}{\varepsilon}
\newcommand{\dd}{\,\mathrm{d}}
\begin{document}
\section*{UnsolvedMath ID: 30002581}
\textbf{Source problem code:} catalogue entry
\subsection{Definitions and mathematical statement}
For reference, under the standard one-variable Cowen--Douglas convention, let \(\Omega\subseteq\C\) be a nonempty connected open set and let \(\mathcal H\) be a complex Hilbert space. For a positive integer \(m\), the class \(B_m(\Omega)\) consists of bounded linear operators \(T\colon\mathcal H\to\mathcal H\) such that, for every \(w\in\Omega\),
\[
\operatorname{ran}(T-wI)=\mathcal H,
\qquad \dim\ker(T-wI)=m,
\qquad
\overline{\operatorname{span}}\!\left(\bigcup_{w\in\Omega}\ker(T-wI)\right)=\mathcal H.
\]
The source does not specify the domain or its definition of \(B_m(\Omega)\); this convention makes the notation explicit but requires confirmation against the cited report.

The supplied rank-two formula is
\[
FB_2(\Omega)=\left\{
T=\begin{pmatrix}T_0&S\\0&T_1\end{pmatrix}
\text{ on }\mathcal H_0\oplus\mathcal H_1:
T_i\in B_1(\Omega)\ (i=0,1),\quad
S\colon\mathcal H_1\to\mathcal H_0\text{ bounded and linear},\quad
T_0S=ST_1
\right\}.
\]
This displayed definition permits \(S=0\). No nonvanishing condition is supplied, so none is added here; whether such a condition is intended is unresolved.

\textbf{Research problem, conditional on the intended flag-class definition.} For each integer \(n>2\), develop a natural higher-rank Cowen--Douglas flag class \(FB_n(\Omega)\) extending the displayed rank-two class, describe its structure, determine complete unitary invariants, and give a reproducing-kernel operator model. Here a complete unitary invariant means data \(\mathcal I(T)\) satisfying
\[
\mathcal I(T)=\mathcal I(\widetilde T)
\quad\Longleftrightarrow\quad
\text{there exists a unitary }U\text{ with }UT=\widetilde T U.
\]
A reproducing-kernel operator model should specify the function space, its kernel, and an operator on that space unitarily equivalent to \(T\), including any necessary conjugation of the parameter domain.

The original wording asks for the structure and model of an apparently already designated class \(FB_n(\Omega)\), whereas the supplied edited wording asks to develop such a class. No higher-rank definition, flag conditions, meaning of ``natural,'' or required model normalization is supplied. Thus the record determines a research objective but does not determine a unique higher-rank classification statement. If a flag is additional data, whether equivalence must also preserve that flag requires clarification.
\subsection{Short English statement}
For every integer \(n>2\), describe the intended higher-rank Cowen--Douglas flag class extending the supplied rank-two formula, classify its operators up to unitary equivalence, and construct reproducing-kernel models. The higher-rank class and some rank-two hypotheses remain unspecified, so this is a conditional statement-only draft requiring mathematical review.
\subsection{Sources}
\begin{itemize}
\item Supplied local UnsolvedMath registry record, permanent ID \(30002581\), titled \emph{Operator Models for Higher-Rank Cowen--Douglas Flag Classes}; exact source problem code \texttt{catalogue entry}, in the Oberwolfach Reports Open Problems collection. The record gives proposed year 2014 and names no proposer.
\item Supplied report citation: \emph{Hilbert Modules and Complex Geometry (2015)}, Oberwolfach Reports, pp.~1139--1219, DOI \texttt{10.4171/owr/2014/21}: \url{https://doi.org/10.4171/owr/2014/21}. The supplied title year is retained alongside the DOI containing 2014; the bibliographic year requires confirmation. No problem-specific page is supplied.
\item Supplied literature source, ``Flag-structure classification'': \url{https://doi.org/10.1016/j.crma.2014.04.001}.
\item Supplied literature source, ``Strong flag structure'': \url{https://arxiv.org/abs/2311.15491}.
\item Supplied literature source, ``Oberwolfach report'': \url{https://doi.org/10.4171/owr/2014/21}.
\item Literature assessment dated exactly \texttt{2026-08-21}: the supplied curation labels the problem partially solved and reports complete unitary invariants for the rank-two flag class, a later curvature-data classification for a strongly flag subclass, and no general answer for all operators in \(FB_n(\Omega)\). These are supplied dated curation claims, not independently verified mathematical or literature conclusions. The background expressly calls for independent verification of time-sensitive status and cited claims.
\item The registry labels the statement \texttt{corrected\_verified}, explaining that mathematical content was preserved while notation, typography, or omitted definitions were repaired for publication. This supplied editorial label does not establish mathematical correctness or resolve the formulation issues recorded here.
\end{itemize}
\end{document}
